\documentclass[a4paper]{article}

\usepackage[margin=2cm]{geometry}
\usepackage{fontspec}
\usepackage[catalan]{babel}
\usepackage{graphicx}

\title{Construcció d'una orquestra amb Raspberry Pi}
\author{Arnau Canyadell Miquel \and Joan Marcè Igual}
\date{Primavera de 2017}

\begin{document}

\begin{titlepage}
	\centering
	\vspace{1cm}
	\includegraphics[width=0.25\textwidth]{images/etseib}
	\par\vspace{1cm}
	\textsc{ \LARGE Escola Tècnica Superior d'Enginyeria \\[1em] 
		Industrial de Barcelona}
	\par\vspace{2cm}
	\textbf{\Huge Projecte II}
	\par\vspace{2cm}
	{\LARGE Construcció d'una orquestra amb Raspberry Pi}
	\vfill
	\begin{flushright}
		\large
		Arnau Canyadell i Miquel \par
		Joan Marcè i Igual \par
	\end{flushright}
\end{titlepage}

\tableofcontents

\newpage

\section{Introducció}

\subsection{Motivacions}

Donat que els dos creadors d'aquest projecte som músics i que tots dos hem participat en diverses orquestres el projecte de construir una orquestra mitjançant la tecnologia i \emph{Internet of Things} ens va semblar molt interessant i que podria integrar-se bé amb els nostres coneixements.

\subsection{Altres projectes}

Buscant si hi ha altres projectes semblants al que es vol realitzar s'ha trobat que Google l'any 2011 va crear el projecte \emph{YouTube Symphony Orchestra 2011}

\section{Objectius}

\section{Desenvolupament del treball}

\section{Resultats}

\section{Bibliografia}


\end{document}